\documentclass[aspectratio=169]{beamer}
\usepackage[utf8]{inputenc}
\usepackage{xeCJK}
\usepackage{graphicx}
\usepackage{booktabs}
\usepackage{tikz}
\usepackage{pgfplots}
\usepackage{xcolor}
\usepackage{colortbl}
\usepackage{longtable}
\usepackage{array}

% 設定中文字體
\setCJKmainfont{Noto Sans CJK TC}

% 主題設定
\usetheme{Madrid}
\usecolortheme{seahorse}

% 標題頁資訊
\title{胸腔內科可信賴專業活動EPAs}
\subtitle{Pulmonology Entrustable Professional Activities (14 EPAs)}
\author{內科部胸腔科}
\date{\today}
\institute{醫學中心}

% 自訂顏色 - 沿用原文件配色
\definecolor{emergency_red}{RGB}{186,24,27}
\definecolor{emergency_blue}{RGB}{0,114,188}
\definecolor{emergency_gray}{RGB}{120,120,120}
\definecolor{highlight_yellow}{RGB}{255,250,205}

% 調整beamer顏色主題
\setbeamercolor{structure}{fg=emergency_blue}
\setbeamercolor{frametitle}{fg=emergency_red}
\setbeamercolor{block title}{bg=emergency_blue!20,fg=emergency_blue}
\setbeamercolor{block body}{bg=emergency_blue!5}
\setbeamercolor{block title alerted}{bg=emergency_red!20,fg=emergency_red}
\setbeamercolor{block body alerted}{bg=emergency_red!5}

\begin{document}

% 標題頁
\begin{frame}
\titlepage
\end{frame}

% 大綱
\begin{frame}{今日議程}
\tableofcontents
\end{frame}

\section{胸腔內科EPAs概述}

\begin{frame}{什麼是胸腔內科EPAs？}
\begin{block}{Entrustable Professional Activities (EPA)}
\textbf{可信賴專業活動} - 胸腔內科醫師在臨床實務中能夠獨立勝任的專業工作單元
\end{block}

\vspace{0.5cm}

\begin{columns}[c]
\column{0.5\textwidth}
\textbf{涵蓋範圍：}
\begin{itemize}
\item 呼吸系統疾病診斷
\item 急慢性肺病處理
\item 胸腔檢查技術
\item 病患照護協調
\item \textcolor{emergency_red}{結核病診療管理}
\item \textcolor{emergency_red}{肺癌診斷分期}
\item \textcolor{emergency_red}{間質性肺疾病}
\end{itemize}

\column{0.5\textwidth}
\textbf{核心特色：}
\begin{itemize}
\item 以能力為導向
\item 漸進式發展
\item 真實情境評量
\item 個別化學習
\item \textcolor{emergency_red}{長期個案追蹤}
\item \textcolor{emergency_red}{多科團隊合作}
\end{itemize}
\end{columns}
\end{frame}

\begin{frame}{委託等級架構}
\begin{center}
\begin{tikzpicture}
% 定義顏色
\definecolor{level1}{RGB}{255,182,193}
\definecolor{level2}{RGB}{255,218,185}
\definecolor{level3}{RGB}{255,255,224}
\definecolor{level4}{RGB}{144,238,144}
\definecolor{level5}{RGB}{173,216,230}

% 繪製發展階梯
\fill[level1] (0,0) rectangle (2,1);
\node at (1,0.5) {\scriptsize Level 1};
\node at (1,-0.5) {\scriptsize Novice};

\fill[level2] (2,0) rectangle (4,1);
\node at (3,0.5) {\scriptsize Level 2};
\node at (3,-0.5) {\scriptsize Advanced};
\node at (3,-0.7) {\scriptsize Beginner};

\fill[level3] (4,0) rectangle (6,1);
\node at (5,0.5) {\scriptsize Level 3};
\node at (5,-0.5) {\scriptsize Competent};

\fill[level4] (6,0) rectangle (8,1);
\node at (7,0.5) {\scriptsize Level 4};
\node at (7,-0.5) {\scriptsize Proficient};

\fill[level5] (8,0) rectangle (10,1);
\node at (9,0.5) {\scriptsize Level 5};
\node at (9,-0.5) {\scriptsize Expert};

% 邊框
\draw[thick] (0,0) rectangle (2,1);
\draw[thick] (2,0) rectangle (4,1);
\draw[thick] (4,0) rectangle (6,1);
\draw[thick] (6,0) rectangle (8,1);
\draw[thick] (8,0) rectangle (10,1);

% 箭頭
\draw[->, thick] (0,-1) -- (10,-1);
\node at (5,-1.5) {能力發展進程};

% 標示重點
\node[above] at (2,1.2) {\scriptsize 需監督};
\node[above] at (6,1.2) {\scriptsize 獨立作業};
\node[above] at (9,1.2) {\scriptsize 指導他人};
\end{tikzpicture}
\end{center}

\vspace{0.5cm}
\begin{itemize}
\item \textbf{Level 1-2}：在監督下學習與執行
\item \textbf{Level 3}：能獨立勝任標準作業
\item \textbf{Level 4-5}：能處理複雜情況並指導他人
\end{itemize}
\end{frame}

\section{14大EPAs架構}

\begin{frame}{胸腔內科14大EPAs概覽（第一部分）}
\footnotesize
\begin{table}
\centering
\begin{tabular}{|p{1cm}|p{3.5cm}|p{6.5cm}|}
\hline
\rowcolor{emergency_blue!20}
\textbf{EPA} & \textbf{專業活動} & \textbf{核心能力領域} \\
\hline
EPA 1 & 病史詢問與風險評估 & 症狀評估、職業暴露、功能分級 \\
\hline
EPA 2 & 胸部理學檢查 & 視觸叩聽技巧、呼吸音識別、異常評估 \\
\hline
EPA 3 & 胸部影像判讀分析 & X光、CT判讀、病變識別、臨床關聯 \\
\hline
EPA 4 & 肺功能檢查操作 & 肺量計操作、功能評估、結果判讀 \\
\hline
EPA 5 & 急性呼吸衰竭處理 & 快速診斷、嚴重度評估、緊急治療 \\
\hline
EPA 6 & 氣道疾病診療 & 氣喘COPD診斷、嚴重度評估、治療選擇 \\
\hline
EPA 7 & 肺炎感染疾病管理 & 病因識別、抗生素選擇、整合照護 \\
\hline
\end{tabular}
\end{table}
\end{frame}

\begin{frame}{胸腔內科14大EPAs概覽（第二部分）}
\footnotesize
\begin{table}
\centering
\begin{tabular}{|p{1cm}|p{3.5cm}|p{6.5cm}|}
\hline
\rowcolor{emergency_blue!20}
\textbf{EPA} & \textbf{專業活動} & \textbf{核心能力領域} \\
\hline
EPA 8 & 肺栓塞診斷管理 & 臨床評分、診斷策略、抗凝治療 \\
\hline
EPA 9 & 胸腔穿刺術 & 適應症評估、解剖定位、併發症預防 \\
\hline
EPA 10 & 支氣管內視鏡檢查 & 設備操作、解剖識別、治療處置 \\
\hline
\rowcolor{emergency_red!10}
EPA 11 & \textcolor{emergency_red}{結核病診斷治療} & \textcolor{emergency_red}{診斷策略、抗結核治療、抗藥性處理} \\
\hline
\rowcolor{emergency_red!10}
EPA 12 & \textcolor{emergency_red}{NTM感染管理} & \textcolor{emergency_red}{菌種識別、ATS標準、治療決策} \\
\hline
\rowcolor{emergency_red!10}
EPA 13 & \textcolor{emergency_red}{肺癌診斷分期} & \textcolor{emergency_red}{早期偵測、分子檢測、多科整合} \\
\hline
\rowcolor{emergency_red!10}
EPA 14 & \textcolor{emergency_red}{間質性肺疾病} & \textcolor{emergency_red}{HRCT判讀、治療策略、進展評估} \\
\hline
\end{tabular}
\end{table}
\end{frame}

\section{重點EPAs詳述}

\begin{frame}{EPA 1: 胸腔科病史詢問與風險評估}
\begin{alertblock}{核心描述}
能夠獨立進行呼吸系統疾病相關的完整病史詢問，包括症狀評估、危險因子識別、職業史調查，並能初步評估呼吸系統風險。
\end{alertblock}

\begin{columns}[c]
\column{0.5\textwidth}
\textbf{關鍵技能：}
\begin{itemize}
\item 咳嗽咳痰詳細評估
\item mMRC呼吸困難分級
\item 吸菸職業暴露史
\item 環境危險因子調查
\end{itemize}

\column{0.5\textwidth}
\textbf{里程碑重點：}
\begin{itemize}
\item Level 2: 系統性病史詢問
\item Level 3: 風險分層評估
\item Level 4: 複雜情況處理
\item Level 5: 指導他人技巧
\end{itemize}
\end{columns}
\end{frame}

\begin{frame}{EPA 5: 急性呼吸衰竭診斷與初步處理}
\begin{alertblock}{核心描述}
能夠快速識別急性呼吸衰竭，進行適當的診斷性檢查，評估嚴重度分級，並執行初步治療措施，包括氧氣治療與呼吸器支持決策。
\end{alertblock}

\begin{columns}[c]
\column{0.5\textwidth}
\textbf{關鍵技能：}
\begin{itemize}
\item 急性呼吸窘迫識別
\item P/F ratio、SOFA評分
\item 氧氣治療選擇
\item NIPPV適應症評估
\end{itemize}

\column{0.5\textwidth}
\textbf{里程碑重點：}
\begin{itemize}
\item Level 2: 標準流程執行
\item Level 3: 獨立診療決策
\item Level 4: 複雜個案處理
\item Level 5: 專業諮詢提供
\end{itemize}
\end{columns}
\end{frame}

\begin{frame}{EPA 9: 胸腔穿刺術與胸管置放}
\begin{alertblock}{核心描述}
能夠執行胸腔穿刺術與胸管置放術，包括適應症評估、解剖定位、無菌技術、併發症預防，並能處理術後照護與管理。
\end{alertblock}

\begin{columns}[c]
\column{0.5\textwidth}
\textbf{關鍵技能：}
\begin{itemize}
\item 診斷性治療性穿刺
\item 超音波導引技術
\item 無菌技術掌握
\item 併發症處理
\end{itemize}

\column{0.5\textwidth}
\textbf{里程碑重點：}
\begin{itemize}
\item Level 2: 監督下執行
\item Level 3: 獨立穿刺操作
\item Level 4: 複雜技術應用
\item Level 5: 專業技術指導
\end{itemize}
\end{columns}
\end{frame}

\begin{frame}{EPA 11: 結核病診斷與治療管理}
\begin{alertblock}{核心描述}
能夠診斷各型結核病，評估傳染性與嚴重度，執行標準抗結核治療，監測治療反應與副作用，並處理抗藥性結核病及其併發症。
\end{alertblock}

\begin{columns}[c]
\column{0.5\textwidth}
\textbf{關鍵技能：}
\begin{itemize}
\item 痰液檢查、IGRA診斷
\item 第一線抗TB藥物
\item 肝毒性監測
\item MDR-TB處理
\item 感染控制措施
\end{itemize}

\column{0.5\textwidth}
\textbf{里程碑重點：}
\begin{itemize}
\item Level 2: 標準治療執行
\item Level 3: 典型TB管理
\item Level 4: 複雜TB處理
\item Level 5: MDR-TB專家
\end{itemize}
\end{columns}
\end{frame}

\begin{frame}{EPA 12: 非典型分枝桿菌感染診斷與管理}
\begin{alertblock}{核心描述}
能夠診斷非典型分枝桿菌感染，區分致病性與定植，執行適當的抗分枝桿菌治療，並管理治療相關併發症與抗藥性問題。
\end{alertblock}

\begin{columns}[c]
\column{0.5\textwidth}
\textbf{關鍵技能：}
\begin{itemize}
\item MAC菌種識別
\item ATS診斷標準
\item 結節支氣管擴張型
\item 藥敏試驗判讀
\item 外科評估
\end{itemize}

\column{0.5\textwidth}
\textbf{里程碑重點：}
\begin{itemize}
\item Level 2: 致病性判定
\item Level 3: NTM治療計畫
\item Level 4: 內外科整合
\item Level 5: NTM專家
\end{itemize}
\end{columns}
\end{frame}

\begin{frame}{EPA 13: 肺癌診斷與分期評估}
\begin{alertblock}{核心描述}
能夠診斷肺癌，執行完整分期評估，協調多專科團隊照護，參與治療決策討論，並提供病患與家屬適當的衛教與支持。
\end{alertblock}

\begin{columns}[c]
\column{0.5\textwidth}
\textbf{關鍵技能：}
\begin{itemize}
\item 肺癌篩檢偵測
\item TNM分期評估
\item EGFR、ALK檢測
\item 腫瘤團隊討論
\item 症狀緩解照護
\end{itemize}

\column{0.5\textwidth}
\textbf{里程碑重點：}
\begin{itemize}
\item Level 2: 分期檢查執行
\item Level 3: 肺癌分期完成
\item Level 4: 多科治療整合
\item Level 5: 肺癌專家領導
\end{itemize}
\end{columns}
\end{frame}

\begin{frame}{EPA 14: 間質性肺疾病診斷與治療}
\begin{alertblock}{核心描述}
能夠診斷各種間質性肺疾病，執行系統性診斷評估，區分可治療與不可治療病因，制定個別化治療方案，並監測疾病進展。
\end{alertblock}

\begin{columns}[c]
\column{0.5\textwidth}
\textbf{關鍵技能：}
\begin{itemize}
\item HRCT影像判讀
\item 藥物性、職業性ILD
\item 支氣管肺泡沖洗
\item 抗纖維化藥物
\item 肺移植評估
\end{itemize}

\column{0.5\textwidth}
\textbf{里程碑重點：}
\begin{itemize}
\item Level 2: 基本ILD診斷
\item Level 3: ILD治療計畫
\item Level 4: 複雜ILD處理
\item Level 5: ILD專家
\end{itemize}
\end{columns}
\end{frame}

\section{評量方法}

\begin{frame}{EPA評量工具對應（更新版）}
\begin{table}
\centering
\footnotesize
\begin{tabular}{|l|l|l|}
\hline
\rowcolor{emergency_blue!20}
\textbf{EPA類型} & \textbf{主要工具} & \textbf{輔助方法} \\
\hline
技能操作類 & DOPS & 直接觀察評估 \\
(EPA 2, 4, 9, 10) & \scriptsize Direct Observation & 同儕回饋 \\
\hline
臨床推理類 & Mini-CEX & 病例討論 \\
(EPA 1,3,5,6,7,8) & \scriptsize Mini-Clinical Exercise & 模擬情境 \\
\hline
\rowcolor{emergency_red!10}
複雜疾病管理類 & Portfolio-based & 長期個案追蹤 \\
\rowcolor{emergency_red!10}
(EPA 11,12,13,14) & Assessment & 多科團隊評量 \\
\hline
整合照護類 & MSF & 360度回饋 \\
(EPA整合面向) & \scriptsize Multi-Source Feedback & 病患家屬評量 \\
\hline
\end{tabular}
\end{table}

\vspace{0.5cm}

\begin{block}{評量頻率建議（更新版）}
\begin{itemize}
\item \textbf{每月重點評量}：3-4個EPAs（包括1個新增EPA）
\item \textbf{季度全面檢討}：所有14個EPAs進展
\item \textbf{年度能力認證}：達標確認與規劃
\end{itemize}
\end{block}
\end{frame}

\begin{frame}{月度評量主題安排（14個EPAs版本）}
\begin{center}
\begin{tikzpicture}
% 繪製時間軸
\draw[thick, emergency_blue] (0,0) -- (14,0);

% 月份標記
\foreach \x/\month in {0/月1, 2.3/月2, 4.6/月3, 6.9/月4, 9.2/月5, 11.5/月6, 14/月7-12} {
  \draw[emergency_blue] (\x,-0.2) -- (\x,0.2);
  \node[below, emergency_blue] at (\x,-0.3) {\tiny \month};
}

% EPAs安排
\node[above, align=center] at (1.15,0.8) {\tiny EPA 1,2,3\\基礎評估};
\node[above, align=center] at (3.45,0.8) {\tiny EPA 4,5,6\\診斷急症};
\node[above, align=center] at (5.75,0.8) {\tiny EPA 7,8\\感染血管};
\node[above, align=center] at (8.05,0.8) {\tiny EPA 9,10\\進階技術};
\node[above, align=center] at (10.35,0.8) {\tiny EPA 11,12\\分枝桿菌};
\node[above, align=center] at (12.75,0.8) {\tiny EPA 13,14\\腫瘤間質};

% 連接線
\foreach \x in {1.15, 3.45, 5.75, 8.05, 10.35, 12.75} {
  \draw[dashed, emergency_gray] (\x,0.2) -- (\x,0.6);
}
\end{tikzpicture}
\end{center}

\vspace{1cm}

\begin{alertblock}{重點提醒}
新增的4個EPAs（11-14）需要特別長期個案追蹤與多專科合作經驗，建議至少6個月的專門訓練。
\end{alertblock}
\end{frame}

\section{實施策略}

\begin{frame}{CCC評量架構}
\begin{block}{Case-based Clinical Competency 特色}
\begin{itemize}
\item \textbf{真實情境評量}：在實際臨床工作中評估
\item \textbf{多次觀察}：持續追蹤能力發展
\item \textbf{漸進委託}：依能力逐步增加責任
\item \textbf{360度回饋}：多方人員提供回饋
\item \textcolor{emergency_red}{\textbf{長期追蹤}：新增EPAs需要6個月以上個案經驗}
\end{itemize}
\end{block}

\vspace{0.5cm}

\begin{columns}[c]
\column{0.5\textwidth}
\textbf{評量重點：}
\begin{itemize}
\item 安全性第一
\item 能力漸進發展
\item 學習從錯誤改進
\item 個別化指導
\end{itemize}

\column{0.5\textwidth}
\textbf{成功關鍵：}
\begin{itemize}
\item 明確期望設定
\item 即時具體回饋
\item 學習資源支持
\item 持續改善機制
\end{itemize}
\end{columns}
\end{frame}

\begin{frame}{實施時程規劃（14個EPAs版本）}
\begin{center}
\begin{tikzpicture}
% 時間軸
\draw[thick] (0,0) -- (14,0);

% 階段標記
\foreach \x/\phase in {0/準備期, 2.8/試行期, 5.6/基礎實施, 8.4/進階實施, 11.2/全面實施, 14/檢討期} {
  \draw (\x,-0.1) -- (\x,0.1);
  \node[below] at (\x,-0.3) {\tiny \phase};
}

% 活動內容
\node[above, align=center] at (1.4,1) {\tiny 教師訓練\\EPA設計};
\node[above, align=center] at (4.2,1) {\tiny 基礎EPA\\試行};
\node[above, align=center] at (7,1) {\tiny EPA 1-10\\實施};
\node[above, align=center] at (9.8,1) {\tiny EPA 11-14\\導入};
\node[above, align=center] at (12.6,1) {\tiny 全面運作\\持續評量};

% 里程碑
\foreach \x/\milestone in {2.8/基礎完成, 5.6/試行成功, 8.4/進階導入, 11.2/全面實施, 14/持續改善} {
  \fill[emergency_red] (\x,0) circle (0.1);
  \node[rotate=45, anchor=west, emergency_red] at (\x,0.2) {\tiny \milestone};
}
\end{tikzpicture}
\end{center}

\vspace{1cm}

\begin{block}{關鍵里程碑}
\begin{itemize}
\item \textbf{3個月}：完成EPA 1-10設計與教師訓練
\item \textbf{6個月}：基礎EPAs試行並調整系統
\item \textbf{9個月}：導入EPA 11-14進階訓練
\item \textbf{12個月}：全面實施並穩定運作
\item \textbf{15個月}：檢討成效並持續改善
\end{itemize}
\end{block}
\end{frame}

\begin{frame}{學習資源與支持（擴展版）}
\begin{columns}[c]
\column{0.5\textwidth}
\textbf{教學資源：}
\begin{itemize}
\item 標準化病患
\item 模擬教學設備
\item 線上學習平台
\item 胸部影像資料庫
\item 肺功能題庫
\item \textcolor{emergency_red}{TB/NTM個案庫}
\item \textcolor{emergency_red}{肺癌MDT紀錄}
\item \textcolor{emergency_red}{ILD影像集}
\end{itemize}

\textbf{評量工具：}
\begin{itemize}
\item 數位評量平台
\item 行動裝置App
\item 即時回饋系統
\item 學習歷程追蹤
\end{itemize}

\column{0.5\textwidth}
\textbf{支持系統：}
\begin{itemize}
\item Mentor制度
\item 同儕學習小組
\item 跨科際交流
\item 國際經驗分享
\item \textcolor{emergency_red}{長期個案追蹤制}
\item \textcolor{emergency_red}{多專科團隊參與}
\end{itemize}

\textbf{品質保證：}
\begin{itemize}
\item 評量者訓練
\item 標準化流程
\item 定期校正會議
\item 外部品質稽核
\item \textcolor{emergency_red}{專科認證連結}
\end{itemize}
\end{columns}
\end{frame}

\section{預期效益}

\begin{frame}{實施14個EPAs的預期效益}
\begin{block}{教育品質提升}
\begin{itemize}
\item 學習目標更明確具體
\item 評量方式更多元實用
\item 個別化學習支持
\item 能力認證更客觀
\item \textcolor{emergency_red}{涵蓋胸腔科完整疾病光譜}
\end{itemize}
\end{block}

\begin{block}{臨床照護改善}
\begin{itemize}
\item 醫師能力標準化
\item 病人安全保障
\item 醫療品質提升
\item 團隊合作強化
\item \textcolor{emergency_red}{複雜疾病管理能力強化}
\end{itemize}
\end{block}

\begin{alertblock}{長期影響}
建立以能力為導向的醫學教育文化，培養更優秀的胸腔內科專科醫師，能處理從常見疾病到複雜罕見疾病的完整疾病光譜
\end{alertblock}
\end{frame}

\begin{frame}{成功實施的關鍵因子}
\begin{center}
\begin{tikzpicture}
% 中心EPAs圓圈
\node[circle, fill=emergency_blue, text=white, font=\Large\bfseries, minimum size=2.5cm] at (0,0) {14個EPAs\\成功實施};

% 四個關鍵因子
\node[rectangle, fill=emergency_red!20, text=black, minimum width=2.8cm, minimum height=1.2cm] at (-3.5,2) {
\begin{tabular}{c}
\textbf{領導支持} \\
高層承諾 \\
資源投入
\end{tabular}
};

\node[rectangle, fill=emergency_red!20, text=black, minimum width=2.8cm, minimum height=1.2cm] at (3.5,2) {
\begin{tabular}{c}
\textbf{教師準備} \\
培訓充分 \\
共識建立
\end{tabular}
};

\node[rectangle, fill=emergency_red!20, text=black, minimum width=2.8cm, minimum height=1.2cm] at (-3.5,-2) {
\begin{tabular}{c}
\textbf{系統整合} \\
流程標準 \\
技術支持
\end{tabular}
};

\node[rectangle, fill=emergency_red!20, text=black, minimum width=2.8cm, minimum height=1.2cm] at (3.5,-2) {
\begin{tabular}{c}
\textbf{持續改善} \\
定期檢討 \\
動態調整
\end{tabular}
};

% 連接線
\foreach \angle in {45, 135, 225, 315} {
  \draw[thick, emergency_blue] (0,0) -- (\angle:2.8cm);
}
\end{tikzpicture}
\end{center}
\end{frame}

\section{新增EPAs特色}

\begin{frame}{新增4個EPAs的特殊考量}
\begin{alertblock}{EPA 11-14的獨特挑戰}
這四個新增的EPAs代表胸腔內科的進階專業領域，需要特別的實施策略
\end{alertblock}

\begin{columns}[c]
\column{0.5\textwidth}
\textbf{共同特點：}
\begin{itemize}
\item 疾病複雜度高
\item 診療週期長
\item 需多科合作
\item 個案相對少見
\item 需長期追蹤
\end{itemize}

\column{0.5\textwidth}
\textbf{實施策略：}
\begin{itemize}
\item Portfolio評量方式
\item 至少6個月訓練
\item 多專科團隊參與
\item 區域醫院合作
\item 國際指引依據
\end{itemize}
\end{columns}

\vspace{0.5cm}

\begin{block}{學習建議}
\begin{itemize}
\item \textbf{EPA 11-12}：感染科合作，建立分枝桿菌診療中心經驗
\item \textbf{EPA 13}：腫瘤科、胸腔外科合作，參與完整MDT流程
\item \textbf{EPA 14}：風濕免疫科合作，建立ILD多專科診療模式
\end{itemize}
\end{block}
\end{frame}

\section{討論與展望}

\begin{frame}{討論議題}
\begin{enumerate}
\item 現有教學資源是否足夠支持14個EPAs實施？
\item 新增EPAs的評量者訓練具體規劃？
\item 與現行住院醫師評量制度如何整合？
\item 學習者接受度與適應性評估？
\item 跨科EPAs的協調與統一標準？
\item 國際認證與接軌的考量？
\item \textcolor{emergency_red}{長期個案追蹤的可行性與挑戰？}
\item \textcolor{emergency_red}{多專科團隊參與的實務困難？}
\end{enumerate}

\vspace{1cm}
\begin{center}
\Large{\textcolor{emergency_blue}{\textbf{歡迎提問與討論}}}
\end{center}
\end{frame}

\begin{frame}{未來展望}
\begin{alertblock}{短期目標 (1年內)}
\begin{itemize}
\item 完成14個胸腔內科EPAs的標準化設計
\item 建立完整的評量工具與流程
\item 訓練足夠的合格評量者
\item 建置數位化學習與評量平台
\item 建立EPA 11-14的長期個案追蹤制度
\end{itemize}
\end{alertblock}

\begin{block}{中期目標 (2-3年)}
\begin{itemize}
\item 擴展到其他內科次專科
\item 建立跨院校合作網絡
\item 發展創新教學方法
\item 進行教學成效研究
\item 與國際EPAs標準接軌
\end{itemize}
\end{block}

\begin{exampleblock}{長期願景 (5年以上)}
\begin{itemize}
\item 成為國內醫學教育標竿
\item 與國際標準接軌認證
\item 推廣至其他醫學專科
\item 建立持續改善機制
\item 培養胸腔內科完整人才
\end{itemize}
\end{exampleblock}
\end{frame}

\begin{frame}{結語}
\begin{center}
\Large{\textcolor{emergency_red}{\textbf{胸腔內科14個EPAs}}}
\end{center}

\vspace{0.5cm}

\begin{center}
\large{不只是評量工具，更是教育革新的起點}
\end{center}

\vspace{1cm}

\begin{itemize}
\item 從知識傳授到能力培養
\item 從標準化到個別化學習
\item 從單向教學到互動發展
\item 從孤立評估到整合成長
\item \textcolor{emergency_red}{從基礎疾病到複雜疾病光譜}
\item \textcolor{emergency_red}{從個人能力到團隊協作}
\end{itemize}

\vspace{1cm}

\begin{center}
\textbf{\textcolor{emergency_blue}{讓我們攜手建立更優質的胸腔內科教育！}}
\end{center}
\end{frame}

\begin{frame}[plain]
\begin{center}
\Huge{\textcolor{emergency_blue}{\textbf{謝謝聆聽}}}

\vspace{1cm}

\Large{\textcolor{emergency_gray}{Questions \& Discussion}}

\vspace{1cm}

\textcolor{emergency_red}{\textbf{內科部胸腔科}}\\
\textcolor{emergency_blue}{Pulmonology Division, Department of Internal Medicine}\\
\vspace{0.5cm}
\textcolor{emergency_gray}{\textit{14 EPAs for Comprehensive Pulmonary Medicine Training}}
\end{center}
\end{frame}

\end{document}